\documentclass[12pt]{article}
\usepackage[utf8]{inputenc}
\usepackage[spanish]{babel}
\usepackage{amsmath}
\usepackage{amsfonts}
\usepackage{amssymb}
\usepackage[margin=2cm]{geometry}
\geometry{a4paper, margin=1in}

\begin{document}

\begin{center}
    Universidad de Antioquia - Instituto de Física\\
    Mecánica de Medios Continuos\\
    Segundo examen -- 08/09/2026
\end{center}

\textbf{NOTA:} es necesario hacer todos los procedimientos matemáticos y las respuestas deben ser simplificadas y justificadas.

\begin{enumerate}
    \item (40\%) \textbf{Atmósfera de dos capas.} Considere una atmósfera en equilibrio hidrostático sometida a un campo gravitacional constante $\vec{g} = -g_0 \hat{k}$. El aire satisface la ecuación de estado del gas ideal
    \[
    p = \rho R_s T,
    \]
    donde $R_s$ es la constante específica del gas, igual en ambas capas. La atmósfera está compuesta por dos regiones:
    \begin{itemize}
        \item para $0 \le z \le z_t$, el aire satisface la relación politrópica $p = K\rho^\gamma$, con $K > 0$ y $\gamma > 1$ constantes;
        \item para $z > z_t$, la atmósfera es isotérmica, con temperatura $T_t = T(z_t)$.
    \end{itemize}
    En la superficie, $p(0) = p_0$ y $T(0) = T_0$. En la tropopausa, situada en $z = z_t$, la presión, la densidad y la temperatura son continuas.
    \begin{itemize}
        \item[a)] Obtenga los perfiles de presión $p(z)$ y temperatura $T(z)$ en la capa inferior en términos de $p_0$, $T_0$, $R_s$, $g_0$ y $\gamma$. Determine sus valores en la tropopausa, $p_t = p(z_t)$ y $T_t = T(z_t)$.
        \item[b)] Obtenga el perfil de presión $p(z)$ en la región isotérmica, imponiendo la continuidad en $z = z_t$.
        \item[c)] Sustituya los valores de $T_0$, $p_0$, $z_t$, $R_s$, $\gamma$ y $g_0$ en los resultados anteriores y evalúe la presión a una altura $z = 20\,\mathrm{km}$.
    \end{itemize}

\section{Solución al punto 1 del Parcial 2: Atmósfera de dos capas}

Se considera una atmósfera en equilibrio hidrostático bajo el campo gravitacional
$\vec{g}=-g_{0}\hat{k}$, con ecuación de estado del gas ideal $p=\rho R_{s}T$
y $R_{s}$ igual en ambas capas. La ecuación de equilibrio hidrostático es
$$\frac{dp}{dz}=-\rho g_{0}.$$

\subsection{Ítem a: perfiles en la capa inferior $(0\le z\le z_{t})$}

En la capa inferior la relación es politrópica, $p=K\rho^{\gamma}$ con $K>0$ y $\gamma>1$.
De aquí se despeja la densidad:
$$\rho=\left(\frac{p}{K}\right)^{1/\gamma}.$$

Sustituyendo en la ecuación hidrostática:
$$\frac{dp}{dz}=-g_{0}\left(\frac{p}{K}\right)^{1/\gamma}
=-\frac{g_{0}}{K^{1/\gamma}}\,p^{1/\gamma}.$$

Separando variables,
$$p^{-1/\gamma}\,dp=-\frac{g_{0}}{K^{1/\gamma}}\,dz.$$

Integrando, usando $\displaystyle \int p^{-1/\gamma}dp=\frac{\gamma}{\gamma-1}p^{(\gamma-1)/\gamma}$:
$$\frac{\gamma}{\gamma-1}p^{(\gamma-1)/\gamma}
=-\frac{g_{0}}{K^{1/\gamma}}z+C.$$

La condición de frontera en la superficie es $p(0)=p_{0}$, de modo que
$$C=\frac{\gamma}{\gamma-1}p_{0}^{(\gamma-1)/\gamma}.$$

Por tanto,
$$p^{(\gamma-1)/\gamma}
=p_{0}^{(\gamma-1)/\gamma}-\frac{\gamma-1}{\gamma}\frac{g_{0}}{K^{1/\gamma}}z.$$

Para expresar el resultado en términos de $p_{0},T_{0},R_{s},g_{0},\gamma$, se usa que en la superficie
$p_{0}=K\rho_{0}^{\gamma}$ con $\rho_{0}=p_{0}/(R_{s}T_{0})$. Entonces
$$K^{1/\gamma}=\frac{p_{0}^{1/\gamma}}{\rho_{0}}
=\frac{p_{0}^{1/\gamma}R_{s}T_{0}}{p_{0}}
=\frac{R_{s}T_{0}}{p_{0}^{(\gamma-1)/\gamma}},$$
y así
$$\frac{1}{K^{1/\gamma}}=\frac{p_{0}^{(\gamma-1)/\gamma}}{R_{s}T_{0}}.$$

Sustituyendo:
$$p^{(\gamma-1)/\gamma}
=p_{0}^{(\gamma-1)/\gamma}\left[1-\frac{\gamma-1}{\gamma}\frac{g_{0}z}{R_{s}T_{0}}\right].$$

Elevando a la potencia $\gamma/(\gamma-1)$, el perfil de presión queda
$$\boxed{\,p(z)=p_{0}\left[1-\frac{\gamma-1}{\gamma}\frac{g_{0}z}{R_{s}T_{0}}\right]^{\gamma/(\gamma-1)},\qquad 0\le z\le z_{t}.\,}$$

Es útil definir la escala de altura politrópica
$$h_{2}=\frac{\gamma}{\gamma-1}\frac{R_{s}T_{0}}{g_{0}},$$
de modo que $p(z)=p_{0}(1-z/h_{2})^{\gamma/(\gamma-1)}$.

Para la temperatura, se usa $T=p/(\rho R_{s})$. Como $\rho=(p/K)^{1/\gamma}$,
$$T=\frac{p}{R_{s}}\left(\frac{K}{p}\right)^{1/\gamma}
=\frac{K^{1/\gamma}}{R_{s}}p^{(\gamma-1)/\gamma}.$$

Dividiendo entre el valor superficial $T_{0}=K^{1/\gamma}p_{0}^{(\gamma-1)/\gamma}/R_{s}$:
$$\frac{T(z)}{T_{0}}=\frac{p(z)^{(\gamma-1)/\gamma}}{p_{0}^{(\gamma-1)/\gamma}}
=1-\frac{\gamma-1}{\gamma}\frac{g_{0}z}{R_{s}T_{0}}.$$

Por tanto, el perfil de temperatura es lineal:
$$\boxed{\,T(z)=T_{0}\left[1-\frac{\gamma-1}{\gamma}\frac{g_{0}z}{R_{s}T_{0}}\right]
=T_{0}-\frac{\gamma-1}{\gamma}\frac{g_{0}}{R_{s}}z,\qquad 0\le z\le z_{t}.\,}$$

El gradiente térmico es constante:
$$\frac{dT}{dz}=-\frac{\gamma-1}{\gamma}\frac{g_{0}}{R_{s}}=-\frac{g_{0}}{c_{p}},
\qquad c_{p}=\frac{\gamma R_{s}}{\gamma-1},$$
que corresponde al lapse rate adiabático seco.

Valores en la tropopausa $z=z_{t}$:
$$\boxed{\,T_{t}=T(z_{t})=T_{0}\left[1-\frac{\gamma-1}{\gamma}\frac{g_{0}z_{t}}{R_{s}T_{0}}\right]\,}$$
y
$$\boxed{\,p_{t}=p(z_{t})
=p_{0}\left[1-\frac{\gamma-1}{\gamma}\frac{g_{0}z_{t}}{R_{s}T_{0}}\right]^{\gamma/(\gamma-1)}
=p_{0}\left(\frac{T_{t}}{T_{0}}\right)^{\gamma/(\gamma-1)}.\,}$$

La densidad en la tropopausa, por continuidad y gas ideal, es
$$\rho_{t}=\frac{p_{t}}{R_{s}T_{t}}=\rho_{0}\left(\frac{T_{t}}{T_{0}}\right)^{1/(\gamma-1)}.$$

\subsection{Ítem b: perfil en la región isotérmica $(z>z_{t})$}

Para $z>z_{t}$ la atmósfera es isotérmica con $T=T_{t}=T(z_{t})$ constante.
Con el gas ideal, $\rho=p/(R_{s}T_{t})$. La ecuación hidrostática queda
$$\frac{dp}{dz}=-\frac{p}{R_{s}T_{t}}g_{0}.$$

Separando variables:
$$\frac{dp}{p}=-\frac{g_{0}}{R_{s}T_{t}}\,dz.$$

Integrando desde la tropopausa $z=z_{t}$, donde por continuidad $p(z_{t})=p_{t}$, hasta una altura $z>z_{t}$:
$$\int_{p_{t}}^{p(z)}\frac{dp'}{p'}
=-\frac{g_{0}}{R_{s}T_{t}}\int_{z_{t}}^{z}dz',$$
$$\ln\frac{p(z)}{p_{t}}=-\frac{g_{0}(z-z_{t})}{R_{s}T_{t}}.$$

Por tanto, el perfil de presión en la capa isotérmica es
$$\boxed{\,p(z)=p_{t}\exp\left[-\frac{g_{0}(z-z_{t})}{R_{s}T_{t}}\right],
\qquad z\ge z_{t},\,}$$
con escala de altura isotérmica
$$H_{t}=\frac{R_{s}T_{t}}{g_{0}}.$$

Sustituyendo $p_{t}$ del ítem a, el perfil completo queda expresado en términos de los datos:
$$p(z)=p_{0}\left[1-\frac{\gamma-1}{\gamma}\frac{g_{0}z_{t}}{R_{s}T_{0}}\right]^{\gamma/(\gamma-1)}
\exp\left[-\frac{g_{0}(z-z_{t})}{R_{s}T_{0}\left(1-\dfrac{\gamma-1}{\gamma}\dfrac{g_{0}z_{t}}{R_{s}T_{0}}\right)}\right].$$

En esta capa la densidad y la temperatura satisfacen
$$\rho(z)=\frac{p(z)}{R_{s}T_{t}}
=\rho_{t}\exp\left[-\frac{g_{0}(z-z_{t})}{R_{s}T_{t}}\right],
\qquad T(z)=T_{t},\qquad z\ge z_{t}.$$

La continuidad de presión, densidad y temperatura en $z=z_{t}$ queda garantizada:
$$p(z_{t}^{-})=p(z_{t}^{+})=p_{t},\qquad
\rho(z_{t}^{-})=\rho(z_{t}^{+})=\rho_{t},\qquad
T(z_{t}^{-})=T(z_{t}^{+})=T_{t}.$$



    \item (30\%) \textbf{Esfera autogravitante.} Considere una esfera de fluido en equilibrio hidrostático bajo su propia gravedad, con radio $R$ y distribución de densidad
    \[
    \rho(r) =
    \begin{cases}
        \rho_c\left(1 - \dfrac{r}{R}\right), & 0 \le r \le R,\\
        0, & r > R,
    \end{cases}
    \]
    donde $\rho_c$ es la densidad central. Suponga que la presión exterior es nula.
    \begin{itemize}
        \item[a)] Determine la distribución de presión $p(r)$ en el interior de la esfera, imponiendo la condición de frontera $p(R) = 0$.
        \item[b)] Calcule la presión central $p_c = p(0)$.
    \end{itemize}

\section{Solución al punto 2 del Parcial 2: Esfera autogravitante}

\subsection{Ítem a: Distribución de presión $p(r)$ en el interior de la esfera}

La condición de equilibrio hidrostático para un fluido bajo su propia gravedad está dada por la ecuación:
$$
\frac{dp}{dr} = -\rho(r) g(r)
$$
donde $g(r)$ es la magnitud de la aceleración gravitacional a una distancia $r$ del centro. Por el teorema de la cáscara (o ley de Gauss para la gravedad), $g(r)$ depende únicamente de la masa encerrada dentro de la esfera de radio $r$, denotada como $M(r)$:
$$
g(r) = \frac{G M(r)}{r^2}
$$

Primero, calculamos la masa encerrada $M(r)$ integrando la distribución de densidad $\rho(r') = \rho_c \left(1 - \frac{r'}{R}\right)$ desde el centro hasta un radio $r$:
$$
M(r) = \int_0^r 4\pi r'^2 \rho(r') \, dr' = 4\pi \rho_c \int_0^r \left(r'^2 - \frac{r'^3}{R}\right) dr'
$$
$$
M(r) = 4\pi \rho_c \left[ \frac{r'^3}{3} - \frac{r'^4}{4R} \right]_0^r = 4\pi \rho_c \left( \frac{r^3}{3} - \frac{r^4}{4R} \right)
$$

Sustituyendo $M(r)$ en la expresión para el campo gravitacional $g(r)$:
$$
g(r) = \frac{G}{r^2} \cdot 4\pi \rho_c \left( \frac{r^3}{3} - \frac{r^4}{4R} \right) = 4\pi G \rho_c \left( \frac{r}{3} - \frac{r^2}{4R} \right)
$$

Ahora, reemplazamos $\rho(r)$ y $g(r)$ en la ecuación de equilibrio hidrostático:
$$
\frac{dp}{dr} = -\rho_c \left(1 - \frac{r}{R}\right) \cdot 4\pi G \rho_c \left( \frac{r}{3} - \frac{r^2}{4R} \right)
$$
$$
\frac{dp}{dr} = -4\pi G \rho_c^2 \left(1 - \frac{r}{R}\right) \left( \frac{r}{3} - \frac{r^2}{4R} \right)
$$

Expandiendo el producto de los polinomios:
$$
\left(1 - \frac{r}{R}\right) \left( \frac{r}{3} - \frac{r^2}{4R} \right) = \frac{r}{3} - \frac{r^2}{4R} - \frac{r^2}{3R} + \frac{r^3}{4R^2} = \frac{r}{3} - \frac{7r^2}{12R} + \frac{r^3}{4R^2}
$$

Por lo tanto, la ecuación diferencial para la presión es:
$$
\frac{dp}{dr} = -4\pi G \rho_c^2 \left( \frac{r}{3} - \frac{7r^2}{12R} + \frac{r^3}{4R^2} \right)
$$

Integramos respecto a $r$ para obtener $p(r)$:
$$
p(r) = -4\pi G \rho_c^2 \int \left( \frac{r}{3} - \frac{7r^2}{12R} + \frac{r^3}{4R^2} \right) dr
$$
$$
p(r) = -4\pi G \rho_c^2 \left( \frac{r^2}{6} - \frac{7r^3}{36R} + \frac{r^4}{16R^2} \right) + C
$$

Para determinar la constante de integración $C$, imponemos la condición de frontera de que la presión exterior es nula, es decir, $p(R) = 0$:
$$
p(R) = -4\pi G \rho_c^2 R^2 \left( \frac{1}{6} - \frac{7}{36} + \frac{1}{16} \right) + C = 0
$$

Calculamos la suma de las fracciones buscando un denominador común, que es 144:
$$
\frac{1}{6} - \frac{7}{36} + \frac{1}{16} = \frac{24}{144} - \frac{28}{144} + \frac{9}{144} = \frac{5}{144}
$$

Despejando $C$:
$$
C = 4\pi G \rho_c^2 R^2 \left( \frac{5}{144} \right) = \frac{5\pi G \rho_c^2 R^2}{36}
$$

Sustituyendo $C$ en la expresión general y factorizando $\frac{\pi G \rho_c^2 R^2}{36}$, obtenemos la distribución de presión final:
$$
p(r) = \frac{\pi G \rho_c^2 R^2}{36} \left[ 5 - 24\left(\frac{r}{R}\right)^2 + 28\left(\frac{r}{R}\right)^3 - 9\left(\frac{r}{R}\right)^4 \right]
$$

\subsection{Ítem b: Presión central $p_c$}

La presión central $p_c$ corresponde a la evaluación de la distribución de presión en el origen, es decir, cuando $r = 0$. 

Sustituyendo $r = 0$ en la expresión derivada en el ítem anterior:
$$
p_c = p(0) = \frac{\pi G \rho_c^2 R^2}{36} \left[ 5 - 24(0)^2 + 28(0)^3 - 9(0)^4 \right]
$$

Por lo tanto, la presión en el centro de la esfera autogravitante es:
$$
p_c = \frac{5\pi G \rho_c^2 R^2}{36}
$$


    \item (30\%) \textbf{Esfera politrópica.} Considere una esfera de fluido autogravitante en equilibrio hidrostático, descrita por la ecuación de estado
    \[
    p = K\rho^2,
    \]
    con $K > 0$ constante. Introduzca las variables
    \[
    \rho(r) = \rho_c\theta(\xi),
    \qquad r = a\xi,
    \qquad a = \left(\frac{K}{2\pi G}\right)^{1/2},
    \]
    donde $\rho_c$ es la densidad central y $G$ es la constante de gravitación universal.
    \begin{itemize}
        \item[a)] A partir de las ecuaciones de equilibrio hidrostático y de conservación de masa, demuestre que
        \[
        \frac{1}{\xi^2}\frac{d}{d\xi}
        \left(\xi^2\frac{d\theta}{d\xi}\right) = -\theta.
        \]
        \item[b)] Resuelva esta ecuación suponiendo una solución regular en el centro y normalizada mediante $\theta(0) = 1$.
        \item[c)] Determine el radio de la esfera, definido por el primer cero positivo de $\theta$, el perfil de densidad $\rho(r)$ y la masa total.
    \end{itemize}
\end{enumerate}

\section{Punto 3: Esfera politrópica}

Se considera una esfera de fluido autogravitante en equilibrio hidrostático, con ecuación de estado
$$
p=K\rho^{2}, \qquad K>0,
$$
y se introducen las variables
$$
\rho(r)=\rho_c\,\theta(\xi), \qquad r=a\xi, \qquad 
a=\left(\frac{K}{2\pi G}\right)^{1/2},
$$
donde $\rho_c$ es la densidad central.

\subsection{Ítem a: Obtención de la ecuación adimensional}

Las ecuaciones de equilibrio hidrostático y de conservación de masa para una esfera simétrica son
$$
\frac{dp}{dr}=-\rho\,\frac{Gm(r)}{r^{2}},
$$
$$
\frac{dm}{dr}=4\pi r^{2}\rho .
$$

De la ecuación hidrostática,
$$
\frac{r^{2}}{\rho}\frac{dp}{dr}=-Gm(r).
$$
Derivando respecto de $r$,
$$
\frac{d}{dr}\left(\frac{r^{2}}{\rho}\frac{dp}{dr}\right)
=-G\frac{dm}{dr}.
$$
Usando la conservación de masa,
$$
\frac{d}{dr}\left(\frac{r^{2}}{\rho}\frac{dp}{dr}\right)
=-4\pi G r^{2}\rho .
$$

Como $p=K\rho^{2}$, se tiene
$$
\frac{dp}{dr}=2K\rho\frac{d\rho}{dr},
$$
y por tanto
$$
\frac{1}{\rho}\frac{dp}{dr}=2K\frac{d\rho}{dr}.
$$
Entonces,
$$
\frac{d}{dr}\left(2K r^{2}\frac{d\rho}{dr}\right)
=-4\pi G r^{2}\rho .
$$
Como $K$ es constante,
$$
\frac{1}{r^{2}}\frac{d}{dr}\left(r^{2}\frac{d\rho}{dr}\right)
=-\frac{2\pi G}{K}\rho .
$$

Ahora se introducen las variables adimensionales
$$
\rho(r)=\rho_c\theta(\xi), \qquad r=a\xi .
$$
Se tiene
$$
\frac{d}{dr}=\frac{1}{a}\frac{d}{d\xi},
$$
y
$$
\frac{d\rho}{dr}=\frac{\rho_c}{a}\frac{d\theta}{d\xi}.
$$
Por tanto,
$$
r^{2}\frac{d\rho}{dr}
=a^{2}\xi^{2}\frac{\rho_c}{a}\frac{d\theta}{d\xi}
=a\rho_c\xi^{2}\theta',
$$
donde $\theta'=d\theta/d\xi$. Luego,
$$
\frac{d}{dr}\left(r^{2}\frac{d\rho}{dr}\right)
=\frac{1}{a}\frac{d}{d\xi}
\left(a\rho_c\xi^{2}\theta'\right)
=\rho_c\frac{d}{d\xi}\left(\xi^{2}\theta'\right).
$$
Así,
$$
\frac{1}{r^{2}}\frac{d}{dr}\left(r^{2}\frac{d\rho}{dr}\right)
=\frac{1}{a^{2}\xi^{2}}\rho_c
\frac{d}{d\xi}\left(\xi^{2}\theta'\right).
$$
Reemplazando en la ecuación previa,
$$
\frac{\rho_c}{a^{2}\xi^{2}}
\frac{d}{d\xi}\left(\xi^{2}\theta'\right)
=-\frac{2\pi G}{K}\rho_c\theta .
$$
Cancelando $\rho_c$ y multiplicando por $a^{2}$,
$$
\frac{1}{\xi^{2}}
\frac{d}{d\xi}\left(\xi^{2}\frac{d\theta}{d\xi}\right)
=-\frac{2\pi G a^{2}}{K}\theta .
$$
Como
$$
a^{2}=\frac{K}{2\pi G},
$$
se obtiene finalmente
$$
\boxed{
\frac{1}{\xi^{2}}
\frac{d}{d\xi}\left(\xi^{2}\frac{d\theta}{d\xi}\right)
=-\theta
}.
$$

\subsection{Ítem b: Solución regular con $\theta(0)=1$}

La ecuación obtenida puede escribirse como
$$
\theta''+\frac{2}{\xi}\theta'+\theta=0,
$$
donde las primas indican derivadas respecto de $\xi$.

Para resolverla, se introduce la variable auxiliar
$$
u(\xi)=\xi\theta(\xi).
$$
Entonces,
$$
\theta=\frac{u}{\xi}.
$$
Derivando,
$$
\theta'=\frac{u'}{\xi}-\frac{u}{\xi^{2}},
$$
$$
\theta''=\frac{u''}{\xi}-\frac{2u'}{\xi^{2}}+\frac{2u}{\xi^{3}}.
$$
Sustituyendo en la ecuación,
$$
\left(\frac{u''}{\xi}-\frac{2u'}{\xi^{2}}+\frac{2u}{\xi^{3}}\right)
+\frac{2}{\xi}
\left(\frac{u'}{\xi}-\frac{u}{\xi^{2}}\right)
+\frac{u}{\xi}=0.
$$
Los términos con $u'$ y $u/\xi^{3}$ se cancelan, quedando
$$
\frac{u''}{\xi}+\frac{u}{\xi}=0.
$$
Para $\xi\neq 0$,
$$
u''+u=0.
$$
La solución general es
$$
u(\xi)=A\sin\xi+B\cos\xi.
$$
Por tanto,
$$
\theta(\xi)=\frac{A\sin\xi+B\cos\xi}{\xi}.
$$

La solución debe ser regular en el centro, $\xi=0$. Como
$$
\frac{\cos\xi}{\xi}\sim \frac{1}{\xi}
\qquad (\xi\to 0),
$$
el término con $B$ diverge. Para evitar la divergencia se debe tomar
$$
B=0.
$$
Así,
$$
\theta(\xi)=A\frac{\sin\xi}{\xi}.
$$
Usando el límite
$$
\lim_{\xi\to 0}\frac{\sin\xi}{\xi}=1,
$$
la condición de normalización $\theta(0)=1$ implica
$$
A=1.
$$
Por tanto, la solución regular es
$$
\boxed{
\theta(\xi)=\frac{\sin\xi}{\xi}
}.
$$
En el centro se entiende que
$$
\theta(0)=\lim_{\xi\to 0}\frac{\sin\xi}{\xi}=1.
$$
Además,
$$
\theta'(0)=0,
$$
como corresponde a una solución regular y simétrica en el centro de la esfera.

\subsection{Ítem c: Radio, perfil de densidad y masa total}

El radio de la esfera se define como el primer cero positivo de $\theta(\xi)$. Como
$$
\theta(\xi)=\frac{\sin\xi}{\xi},
$$
los ceros están determinados por
$$
\sin\xi=0.
$$
El primer cero positivo es
$$
\xi_1=\pi.
$$
Por tanto, el radio físico de la esfera es
$$
R=a\xi_1=a\pi.
$$
Usando
$$
a=\left(\frac{K}{2\pi G}\right)^{1/2},
$$
se obtiene
$$
\boxed{
R=\pi\left(\frac{K}{2\pi G}\right)^{1/2}
}.
$$
Equivalentemente,
$$
\boxed{
R=\left(\frac{\pi K}{2G}\right)^{1/2}
}.
$$

El perfil de densidad está dado por
$$
\rho(r)=\rho_c\theta(\xi)
=\rho_c\frac{\sin\xi}{\xi}.
$$
Como $\xi=r/a$ y $R=\pi a$, también se puede escribir
$$
\xi=\frac{r}{a}=\frac{\pi r}{R}.
$$
Por tanto,
$$
\boxed{
\rho(r)=\rho_c\,
\frac{\sin\left(\dfrac{\pi r}{R}\right)}
{\dfrac{\pi r}{R}},
\qquad 0\le r\le R
}.
$$
En el centro se toma el límite
$$
\rho(0)=\rho_c.
$$
Fuera de la esfera, la densidad es nula:
$$
\boxed{
\rho(r)=0,\qquad r>R
}.
$$

La masa total de la esfera es
$$
M=\int_0^R 4\pi r^{2}\rho(r)\,dr.
$$
Usando $r=a\xi$ y $dr=a\,d\xi$,
$$
M=4\pi a^{3}\rho_c\int_0^{\pi}\xi^{2}\theta(\xi)\,d\xi.
$$
Como $\theta(\xi)=\sin\xi/\xi$,
$$
M=4\pi a^{3}\rho_c\int_0^{\pi}\xi\sin\xi\,d\xi.
$$
La integral se calcula por partes:
$$
\int \xi\sin\xi\,d\xi
=-\xi\cos\xi+\sin\xi.
$$
Entonces,
$$
\int_0^{\pi}\xi\sin\xi\,d\xi
=\left[-\xi\cos\xi+\sin\xi\right]_0^{\pi}.
$$
Evaluando,
$$
\left[-\xi\cos\xi+\sin\xi\right]_0^{\pi}
=-\pi\cos\pi+\sin\pi-0
=\pi.
$$
Por tanto,
$$
M=4\pi a^{3}\rho_c\,\pi
=4\pi^{2}a^{3}\rho_c.
$$
Así,
$$
\boxed{
M=4\pi^{2}\rho_c a^{3}
}.
$$
Sustituyendo
$$
a=\left(\frac{K}{2\pi G}\right)^{1/2},
$$
se obtiene
$$
M=4\pi^{2}\rho_c
\left(\frac{K}{2\pi G}\right)^{3/2}.
$$
Equivalentemente,
$$
\boxed{
M=\sqrt{2\pi}\,\rho_c\left(\frac{K}{G}\right)^{3/2}
}.
$$
También puede escribirse en términos del radio $R$. Como $R=\pi a$, entonces $a=R/\pi$, y por tanto
$$
M=4\pi^{2}\rho_c\left(\frac{R}{\pi}\right)^{3}
=\frac{4}{\pi}\rho_c R^{3}.
$$
Así,
$$
\boxed{
M=\frac{4}{\pi}\rho_c R^{3}
}.
$$

En resumen,
$$
\boxed{
\theta(\xi)=\frac{\sin\xi}{\xi}
},
$$
$$
\boxed{
R=\pi\left(\frac{K}{2\pi G}\right)^{1/2}
},
$$
$$
\boxed{
\rho(r)=\rho_c
\frac{\sin\left(\dfrac{\pi r}{R}\right)}
{\dfrac{\pi r}{R}},
\qquad 0\le r\le R
},
$$
$$
\boxed{
M=4\pi^{2}\rho_c a^{3}
=\sqrt{2\pi}\,\rho_c\left(\frac{K}{G}\right)^{3/2}
=\frac{4}{\pi}\rho_c R^{3}
}.
$$


\end{document}
